% Figure: Carnot cycle on the T-s diagram. Two isotherms (horizontal) at T_h
% and T_c, two adiabats (vertical) joining them. The enclosed rectangle is the
% net work; the area below the upper isotherm is the heat absorbed.
\begin{figure}[htbp]
\centering
\begin{tikzpicture}
\begin{axis}[
  width=11cm, height=7.5cm,
  xlabel={specific entropy $s$ (kJ/kg$\cdot$K)},
  ylabel={temperature $T$ (K)},
  xmin=0.5, xmax=4.5, ymin=200, ymax=750,
  axis lines=left,
  xtick={1,4}, xticklabels={$s_1$,$s_2$},
  ytick={300,650}, yticklabels={$T_c$,$T_h$},
  tick label style={font=\scriptsize},
  label style={font=\small},
  clip=false,
]
% rectangle: corners (1,650),(4,650),(4,300),(1,300)
\addplot[engblue, very thick] coordinates {(1,650) (4,650)};
\addplot[engblue, very thick] coordinates {(4,650) (4,300)};
\addplot[engblue, very thick] coordinates {(4,300) (1,300)};
\addplot[engblue, very thick] coordinates {(1,300) (1,650)};
% shade the work rectangle lightly
\addplot[englime!25, fill, draw=none] coordinates {(1,300) (4,300) (4,650) (1,650)} \closedcycle;
% process labels
\node[engblue, font=\scriptsize, anchor=south] at (axis cs:2.5,655) {isothermal expansion, $Q_h$ in};
\node[engblue, font=\scriptsize, anchor=north] at (axis cs:2.5,295) {isothermal compression, $Q_c$ out};
\node[engred, font=\scriptsize, anchor=west, rotate=90] at (axis cs:4.05,475) {adiabat};
\node[engred, font=\scriptsize, anchor=east, rotate=90] at (axis cs:0.95,475) {adiabat};
% state numbers
\node[font=\scriptsize, anchor=south east] at (axis cs:1,650) {1};
\node[font=\scriptsize, anchor=south west] at (axis cs:4,650) {2};
\node[font=\scriptsize, anchor=north west] at (axis cs:4,300) {3};
\node[font=\scriptsize, anchor=north east] at (axis cs:1,300) {4};
\node[font=\scriptsize, anchor=center] at (axis cs:2.5,475) {$W_{\text{net}}$};
\end{axis}
\end{tikzpicture}
\caption[Carnot cycle on the temperature-entropy diagram]{The Carnot cycle on
the temperature-entropy plane is a rectangle. Heat enters along the upper
isotherm at $T_h$ (process 1 to 2) and leaves along the lower isotherm at $T_c$
(process 3 to 4); the two adiabats (2 to 3 and 4 to 1) are vertical because a
reversible adiabatic process holds entropy constant. The enclosed area equals
the net work per cycle, $W_{\text{net}} = (T_h - T_c)(s_2 - s_1)$, and the heat
absorbed is the area under the upper isotherm, $Q_h = T_h(s_2 - s_1)$, so the
efficiency reduces to $1 - T_c/T_h$.}
\label{fig:vol04ch03:carnot-ts}
\end{figure}
